\documentclass[11pt]{article}

\usepackage[margin=1in]{geometry}
\usepackage[T1]{fontenc}
\usepackage[utf8]{inputenc}
\usepackage{amsmath,amssymb,amsthm}
\usepackage{booktabs,tabularx,array}
\usepackage{float}
\usepackage{xcolor}
\usepackage{hyperref}
\usepackage{setspace}
\usepackage{caption}
\usepackage[round]{natbib}

\hypersetup{
  colorlinks=true,
  linkcolor=blue!50!black,
  citecolor=blue!50!black,
  urlcolor=blue!50!black
}
\setstretch{1.04}
\setlength{\parindent}{0pt}
\setlength{\parskip}{0.42em}
\captionsetup{font=small,labelfont=bf}
\newcolumntype{Y}{>{\raggedright\arraybackslash}X}

\newcommand{\E}{\mathbb{E}}
\newcommand{\V}{\mathbb{V}}
\newcommand{\1}{\mathbf{1}}
\newtheorem{proposition}{Proposition}
\newcommand{\NumClusters}{20}
\newcommand{\ClusterSize}{50}
\newcommand{\GoldATE}{0.100}
\newcommand{\ProxyATE}{0.060}
\newcommand{\EndpointGap}{0.040}
\newcommand{\PredictionRMSE}{0.200}
\newcommand{\BudgetExample}{5}
\newcommand{\ConcentratedZeroMove}{55.3\%}
\newcommand{\DiffuseBudgetSD}{0.000}
\newcommand{\ConcentratedBudgetSD}{0.048}
\newcommand{\ConcentratedTailRisk}{5.3\%}
\newcommand{\TailThreshold}{0.050}
\newcommand{\WorstCaseGap}{0.400}
\newcommand{\WorstCaseGapMultiple}{4}
\newcommand{\HajekConcentrated}{45\%}
\newcommand{\HajekOffsetting}{25\%}
\newcommand{\OffsettingTailPeak}{62\%}
\newcommand{\OffsettingTailPeakBudget}{7}
\newcommand{\ConcentratedNormalTail}{29\%}
\newcommand{\OffsettingSignFlip}{10\%}


\title{\textbf{Validation-Swap Distributions}\\[0.35em]
\large Diagnosing Machine-Scored Outcomes in Experiments}
\author{}
\date{August 2026}

\begin{document}
\maketitle

\begin{abstract}
In an experiment with clustered treatment assignment, a machine-learning or
large language model scores every outcome, and human labels are affordable for
only some clusters. Had assignment been individual, fixed scoring error would
be balanced by design and only treatment-dependent error would bias the
contrast; clustering removes that protection, leaving few effective units for
both balancing and validation. Before coding begins, the researcher must decide how many
clusters to validate, and the decision requires knowing what each budget buys.
Two-phase and prediction-powered estimators correct the treatment effect with
the labeled clusters, and power analysis prices a labeling design by the
estimator's asymptotic variance. When scoring error is concentrated in a few
clusters, the variance misprices the risk: the corrected estimate's
distribution across possible label sets is discrete and far from Normal, with
most draws moving the estimate not at all and a few moving it far. We price
the design by that distribution itself. For each budget $k$ we swap human
labels into a set of $k$ clusters, recompute the treatment effect, and
enumerate all sets exactly, which is feasible in benchmarks, pilots, and fully
coded studies used to plan smaller budgets. As $k$ grows, the mean of the
distribution moves on a straight line from the machine-scored estimate to the
fully human-coded one, whatever the error structure, so everything that
distinguishes labeling designs lies in the dispersion and tails. In a worked
example, two scoring systems agree on RMSE, $R^2$, both endpoints, and the
entire mean path, yet when the error sits in two of twenty clusters, labeling
five clusters leaves the estimate unchanged \ConcentratedZeroMove\ of the
time. With a single real validation draw the full distribution is out of
reach, but its variance at every budget remains estimable without bias, and a
negligibility condition says when a Normal approximation can be trusted: only
when no small set of clusters dominates the correction. The distribution
conditions on the observed pool of gold labels; inference about the population
still requires a probability validation design or a transport model.
\end{abstract}

Researchers can score every response in an experiment with a classifier, a
computer-vision model, or a large language model while paying human coders to
score only a fraction \citep{gilardi2023}. The cheap score may have excellent
test-set accuracy and still distort the treatment effect. The relevant error is not average scoring
error. It is the difference in scoring error between treatment arms.

Suppose an experiment tests whether a message changes the hostility expressed in
an open-ended response. A language model scores hostility for all respondents,
and human coders score a random subset. The model can overstate hostility by two
hundredths in one arm and understate it by two hundredths in the other. Its mean
error is zero, but the estimated treatment effect is wrong by four hundredths.
Aggregate accuracy conceals this error because accuracy and the treatment
contrast weight observations differently.

Existing methods already correct downstream inference with a small set of gold
labels. Post-prediction inference, design-based supervised learning, and
prediction-powered inference combine predictions on many observations with
gold-label residuals on a smaller probability sample
\citep{wang2020,egami2023,angelopoulos2023,angelopoulos2023ppi}. Two-phase
designs for allocating expensive measurements go back to
\citet{neyman1938}; modern treatments choose the validation sample to minimize
the variance of a target estimator \citep{amorim2021}, including adaptively
\citep{zrnic2024}. Parallel corrections exist when the machine prediction
enters the analysis as a regression covariate rather than the outcome
\citep{fong2021}. Work on text as data explains why a learned outcome
requires separate measurement and causal design choices \citep{egami2022}.
These methods establish the estimator and its uncertainty.

This paper asks how much the corrected answer depends on which clusters
receive the gold labels, and builds the argument from the case where the
answer is boring. In an individually randomized experiment, scoring error
that is a fixed attribute of the response is balanced by randomization like
any baseline covariate, however strongly it clusters by dialect, site, or
coder batch. Only error that depends on treatment itself biases the contrast,
and an ordinary validation sample corrects that bias with textbook two-phase
theory. Clustered assignment removes this protection twice over: the
balancing now operates over a few clusters rather than many individuals, and
validation samples those same few clusters, so the correction inherits the
randomness of which clusters were checked. We characterize that randomness
exactly through the fixed-budget random-swap distribution of the
treatment-effect estimate. Its spread shows whether a validation budget
produces a stable correction or a lottery; its tails show how often the
estimate moves by more than a prespecified substantive amount.

The contribution is deliberately narrow. The estimator below is a special
case of established difference and prediction-powered estimators, and each
formal result is a short application of standard sampling theory. What is new
is the object: the exact distribution of the corrected estimate over labeling
designs, used to price a validation budget before it is spent.

\section{Machine scores as experimental outcomes}

An experiment measures an outcome for units $i=1,\ldots,N$, of which $N_1$
are treated and $N_0$ are control. Human coding of the outcome $Y_i$ is
expensive, so a machine score $\widehat Y_i$ stands in for it. Define the
scoring error
\begin{equation}
e_i=Y_i-\widehat Y_i.
\label{eq:error}
\end{equation}
The machine score must be fixed independently of the outcome analysis. A score
learned on the experimental outcomes requires a separate training sample or an
appropriate cross-fitting design. Swaps do not repair overfitting.

For clarity, consider the unadjusted difference in means. The proxy and gold
estimates are
\begin{align}
\widehat\tau_P
&=\frac{1}{N_1}\sum_{i:D_i=1}\widehat Y_i
-\frac{1}{N_0}\sum_{i:D_i=0}\widehat Y_i,\\
\widehat\tau_G
&=\frac{1}{N_1}\sum_{i:D_i=1}Y_i
-\frac{1}{N_0}\sum_{i:D_i=0}Y_i,
\end{align}
and their gap is a weighted sum of the errors,
\begin{equation}
\widehat\tau_G-\widehat\tau_P=\sum_{i=1}^{N}w_ie_i,
\qquad
w_i=\frac{D_i}{N_1}-\frac{1-D_i}{N_0}.
\label{eq:gap}
\end{equation}
Whether scoring error is a problem is entirely a question about this weighted
sum. Error orthogonal to the treatment weights leaves the contrast intact no
matter how large the errors are.

Like outcomes, errors have potential values: $e_i(d)=Y_i(d)-\widehat Y_i(d)$
for assignment $d\in\{0,1\}$. Call the error \emph{non-differential} when
$e_i(1)=e_i(0)$, so the score misreads unit $i$'s response by the same amount
under either assignment. \emph{Differential} error, such as a model that
misreads the style treated respondents adopt, violates this condition.

Sections~\ref{sec:base} and~\ref{sec:breaks} take expectations over treatment
assignment to establish orders of magnitude. From
Section~\ref{sec:swap} onward the analysis conditions on the realized
experiment: outcomes, scores, and assignment are fixed, and the only
randomness studied is the validation draw.

\section{The base case: individually randomized experiments}
\label{sec:base}

Assign treatment by complete randomization at the individual level. Under
non-differential error, each $e_i$ is a fixed attribute of unit $i$, and the
gap in equation~\eqref{eq:gap} averages fixed numbers over a random
assignment:
\begin{equation}
\E_R\!\left[\widehat\tau_G-\widehat\tau_P\right]=0,
\qquad
\V_R\!\left[\widehat\tau_G-\widehat\tau_P\right]
=\left(\frac{1}{N_1}+\frac{1}{N_0}\right)S_e^2,
\label{eq:basecase}
\end{equation}
where $S_e^2$ is the finite-population variance of the errors and $\E_R$
averages over assignments. Scoring error behaves exactly like a baseline
covariate: randomization balances it, and its realized imbalance is
$O_p(N^{-1/2})$ noise of the same order as the sampling variability the
standard error already reports.

Group structure in the errors changes nothing. If a group $h$ of $n_h$ units
--- a dialect, a site, a coder batch --- shares a common error component
$\mu_h$, its contribution to the gap is
\begin{equation}
\mu_h\left(\frac{n_{h1}}{N_1}-\frac{n_{h0}}{N_0}\right),
\label{eq:imbalance}
\end{equation}
where $n_{h1}$ and $n_{h0}$ count the group's treated and control units.
Under individual randomization the imbalance in parentheses is
$O_p(\sqrt{n_h}/N)$: group-correlated error cancels across arms because
randomization splits every group.

The one threat that survives the base case is differential error. If
$e_i(1)\neq e_i(0)$, the gap converges to the average error contrast
$\E[e(1)-e(0)]$, a bias no sample size repairs. Even then nothing bespoke is
required: a simple random sample of $m$ units receives human labels, the
paired discrepancies estimate the bias at rate $m^{-1/2}$, and the standard
two-phase and prediction-powered corrections apply with independent-sampling
theory \citep{neyman1938,angelopoulos2023}.

What cannot substitute for those labels is an accuracy audit.

\begin{proposition}[Sharp accuracy bound]
\label{prop:bound}
Let $r=(N^{-1}\sum_ie_i^2)^{1/2}$ be the root mean squared error of the
machine score on the experimental sample. For every error configuration,
\begin{equation}
|\widehat\tau_G-\widehat\tau_P|
\le r\sqrt{N\left(\frac{1}{N_1}+\frac{1}{N_0}\right)},
\label{eq:bound}
\end{equation}
with equality exactly when $e_i$ is proportional to the contrast weight
$w_i$.
\end{proposition}

\begin{proof}
The gap in equation~\eqref{eq:gap} is linear in the errors. By the
Cauchy--Schwarz inequality its absolute value is at most $\|w\|\,\|e\|$,
where $\|w\|^2=1/N_1+1/N_0$ and $\|e\|^2=Nr^2$, with equality exactly when
$e$ is proportional to $w$.
\end{proof}

The bound is attainable, and it is weak at the scale that matters. In the
numerical design below, RMSE \PredictionRMSE\ permits an endpoint gap as
large as \WorstCaseGap, \WorstCaseGapMultiple\ times the gold treatment
effect of \GoldATE. Accuracy metrics constrain the contrast only at the order
of $r$ itself, and they cannot separate differential from non-differential
error. That separation requires labels.

\section{What clustered assignment breaks}
\label{sec:breaks}

Now assign treatment at the cluster level, as in a site-, classroom-, or
wave-randomized experiment: clusters $g=1,\ldots,G$ with members indexed by
$i$, treatment $D_g$ constant within cluster, and errors written $e_{gi}$.
Three things break at once.

First, the balancing in equation~\eqref{eq:basecase} now operates over $G$
clusters rather than $N$ individuals. Over cluster randomization with equal
cluster sizes and equal arms, the gap is still mean zero, but its standard
deviation is $2S_\mu/\sqrt{G}$, where $S_\mu$ is the between-cluster standard
deviation of mean scoring error. With twenty clusters and cluster-level error
variation comparable to the effects under study, the realized experiment
carries a gap of the same order as the estimand. What was vanishing noise is
now a fixed, consequential number.

Second, clusters are maximal-imbalance groups. Each cluster sits wholly in
one arm, so the imbalance factor in equation~\eqref{eq:imbalance} takes its
extreme value for every cluster-shaped error group, and nothing cancels
within arms. Define cluster $g$'s contribution to the endpoint gap as
\begin{equation}
\psi_g=
\begin{cases}
N_1^{-1}\sum_i e_{gi}, & D_g=1,\\
-N_0^{-1}\sum_i e_{gi}, & D_g=0.
\end{cases}
\label{eq:psi}
\end{equation}
The contributions add exactly:
\begin{equation}
\sum_{g=1}^{G}\psi_g=\widehat\tau_G-\widehat\tau_P.
\label{eq:add}
\end{equation}
Equation~\eqref{eq:add} is an accounting identity within the paired
gold-label pool. It does not identify the errors of unlabeled clusters. For
OLS with fixed covariates, residualized-treatment weights replace the treated
and control weights in equation~\eqref{eq:psi}, and the same adding-up and
sampling identities follow; a nonlinear estimator can still be recomputed
over swap sets, but its mean path need not be linear. When scoring failures
are cluster-shaped --- a model that misreads one site's context is wrong for
that whole site --- the gap concentrates in a few large $\psi_g$.

Third, validation must draw whole clusters, because human recoding is
organized in the same batches as the scoring failures. The correction
computed from $k$ labeled clusters is therefore a statistic of $k$ draws from
a population of $G$ numbers that a few of them may dominate: few effective
units, and no central limit comfort when the contributions concentrate. This
is the setting the rest of the paper prices. The same construction applies to
documents, waves, or any other sampling unit, provided the swap unit matches
the validation design and contains the error dependence.

\section{The fixed-budget swap distribution}
\label{sec:swap}

Let $S_k$ be a simple random sample without replacement of exactly $k$ clusters.
The raw-swap estimate replaces machine scores by human outcomes in the selected
clusters:
\begin{equation}
\widehat\tau_R(S_k)=\widehat\tau_P+\sum_{g\in S_k}\psi_g.
\label{eq:raw}
\end{equation}
The distribution of $\widehat\tau_R(S_k)$ over all $\binom{G}{k}$ sets is the
raw fixed-budget swap distribution. It describes partial repair of the observed
data. A fixed $k$ is preferable to independent Bernoulli swaps when $k$
represents a coding budget.

Let $\bar\psi=G^{-1}\sum_g\psi_g$ and
$S_\psi^2=(G-1)^{-1}\sum_g(\psi_g-\bar\psi)^2$. Standard sampling identities
give
\begin{align}
\E_{S_k}[\widehat\tau_R(S_k)]
&=\widehat\tau_P+\frac{k}{G}
(\widehat\tau_G-\widehat\tau_P),
\label{eq:rawmean}\\
\V_{S_k}[\widehat\tau_R(S_k)]
&=k\left(1-\frac{k}{G}\right)S_\psi^2.
\label{eq:rawvar}
\end{align}

Equation~\eqref{eq:rawmean} is the central result. In a linear outcome analysis,
the average raw-swap path must be a straight line connecting the proxy and gold
estimates. Any two scoring systems with the same endpoints have the same mean
path. The path says nothing about whether the correction is spread evenly
across clusters or concentrated in a few. Equation~\eqref{eq:rawvar} shows where
that information resides. Two further facts follow directly from
equations~\eqref{eq:add} and~\eqref{eq:rawmean} and need no proof beyond
arithmetic. The $k=1$ distribution is uniform on
$\{\widehat\tau_P+\psi_g\}$, so it identifies the multiset of cluster
contributions, and every fixed-$k$ distribution is a symmetric function of
that multiset. The mean path, by contrast, depends on the contributions only
through their sum: any two error patterns with the same endpoints share the
entire mean path, however differently the correction is distributed across
clusters.

The distribution supports several descriptive summaries. A researcher can
report its central quantiles, the chance that swapping $k$ clusters changes the
estimate's sign, and the chance that the movement exceeds a prespecified
substantive threshold. These probabilities refer to the stated swap design.
They are not $p$-values.

Raw replacement mechanically shrinks as $k/G$ falls. If the goal is to estimate
the fully human-coded contrast from a $k$-cluster probability sample, scale the
observed correction by its inclusion probability:
\begin{equation}
\widehat\tau_H(S_k)
=\widehat\tau_P+\frac{G}{k}\sum_{g\in S_k}\psi_g.
\label{eq:ht}
\end{equation}
This is a Horvitz-Thompson difference correction
\citep{horvitz1952,sarndal1992}. In this simple setting it is also the
rectifier used by prediction-powered inference \citep{angelopoulos2023}. It
satisfies
\begin{align}
\E_{S_k}[\widehat\tau_H(S_k)]&=\widehat\tau_G,
\label{eq:htmean}\\
\V_{S_k}[\widehat\tau_H(S_k)]
&=\left(\frac{G}{k}\right)^2
\,k\left(1-\frac{k}{G}\right)S_\psi^2.
\label{eq:htvar}
\end{align}
Thus the design-scaled swap distribution is centered on the gold endpoint. Its
spread measures validation-sample fragility at budget $k$.

Unequal-probability validation replaces $G/k$ with inverse inclusion
probabilities. Stratifying the validation draw by treatment arm or by an
observable proxy-risk score can improve precision. Those choices change the
swap distribution and must be included in its definition. The optimal choice
depends on the target estimator \citep{amorim2021,zrnic2024}.

An exact distribution requires $\psi_g$ for every cluster in the finite pool.
That is available in a benchmark dataset, a simulation, or a large pilot used to
study smaller future budgets. With only one small validation sample, the analyst
cannot recover the unseen cluster contributions nonparametrically. Standard
design-based variance estimation remains available when selection probabilities
are known. A cluster bootstrap or a model for $e_{gi}$ can approximate more of
the distribution, but then the result depends on resampling or transport
assumptions.

The second moment of the swap distribution is nevertheless estimable from a
single draw, at the observed budget and at every other one.

\begin{proposition}[The variance path is estimable from one validation sample]
\label{prop:varpath}
Let $S_k$ be a simple random sample of $k\ge 2$ clusters, and let
$s_\psi^2=(k-1)^{-1}\sum_{g\in S_k}(\psi_g-\bar\psi_{S_k})^2$ be the sample
variance of the observed contributions. Then
$\E_{S_k}[s_\psi^2]=S_\psi^2$. Substituting $s_\psi^2$ into
equations~\eqref{eq:rawvar} and~\eqref{eq:htvar} therefore gives unbiased
estimates of the raw and design-scaled swap variances at every budget
$k'\in\{1,\ldots,G\}$, and Chebyshev's inequality bounds the tail probability
$\Pr\{|\widehat\tau_H(S_{k'})-\widehat\tau_G|>\delta\}$ by
$(G/k')^{2}k'(1-k'/G)\,S_\psi^2/\delta^{2}$ for every threshold $\delta>0$.
\end{proposition}

\begin{proof}
Design-unbiasedness of the sample variance under simple random sampling
without replacement is classical \citep{cochran1977,liding2017}. The variance
claims substitute this estimate into the exact formulas, and the tail bound
applies Chebyshev's inequality to the design-scaled estimator, which is
centered on $\widehat\tau_G$ by equation~\eqref{eq:htmean}.
\end{proof}

A pilot at one budget therefore prices fragility at every other budget. This
is the finite-population analogue of power analysis for prediction-powered
inference, which chooses the labeled sample size from an asymptotic variance
\citep{chen2026}. The estimate is honest but noisy: the sampling variance of
$s_\psi^2$ depends on fourth moments of the contributions, so with few
validated clusters and concentrated error the estimated variance path is
itself fragile.

Moment summaries also carry an implicit distributional assumption. The
following negligibility condition, due to \citet{hajek1960} and stated here in
the form of Theorem~1 of \citet{liding2017}, says when the swap distribution
can be treated as Normal.

\begin{proposition}[Normal approximations require diffuse contributions]
\label{prop:normal}
Embed the design in a sequence of finite populations with $G\to\infty$
clusters and budgets $k=k(G)$. If
\begin{equation}
\frac{1}{\min(k,\,G-k)}\cdot
\frac{\max_{g}(\psi_g-\bar\psi)^2}{S_\psi^2}\longrightarrow 0,
\label{eq:hajek}
\end{equation}
then the standardized raw-swap estimate
$\{\widehat\tau_R(S_k)-\E[\widehat\tau_R(S_k)]\}\big/\V[\widehat\tau_R(S_k)]^{1/2}$
converges in distribution to a standard Normal, and the design-scaled
estimate, an affine transformation at fixed $k$, converges with it. The
condition fails when a bounded number of clusters carries a nonvanishing share
of the centered squared contributions.
\end{proposition}

\begin{proof}
The raw-swap estimate is $\widehat\tau_P$ plus $k$ times the sample mean of a
simple random sample of the finite population $\{\psi_1,\ldots,\psi_G\}$, so
the claim is the finite population central limit theorem for simple random
sampling \citep{erdos1959,hajek1960}, in the sufficient form of Theorem~1 of
\citet{liding2017}. Sharp necessary and sufficient conditions of Lindeberg type
are in \citet{hajek1960}.
\end{proof}

The finite-sample diagnostic implied by condition~\eqref{eq:hajek} is the
largest single-cluster share of the centered squared contributions. In the
numerical design below, one cluster carries \HajekConcentrated\ of that total
in the concentrated case and \HajekOffsetting\ in the offsetting case, and the
diffuse case is degenerate. Budget planning that combines
Proposition~\ref{prop:varpath} with a Normal approximation is therefore
reliable exactly when the correction is diffuse, and misleading in the
concentrated case the diagnostic exists to detect. That is the formal reason
to examine the exact distribution rather than its first two moments.

\section{A worked example: same accuracy, different prices}

The numerical example has \NumClusters\ equal-size clusters with
\ClusterSize\ observations each. Half the clusters are treated. The fully
human-coded treatment effect is \GoldATE, and the machine-scored effect is
\ProxyATE\ in the diffuse and concentrated cases. All three scoring systems
have RMSE \PredictionRMSE\ and the same $R^2$ because they score the same gold
outcomes and have the same squared-error total.

The error patterns differ only in their cluster contributions. In the diffuse
case, each cluster contributes $\EndpointGap/\NumClusters$ to the endpoint gap.
In the concentrated case, two clusters contribute $\EndpointGap/2$ each and the
others contribute zero. The offsetting case has two positive and two negative
contributions of equal size, so the proxy and gold endpoints agree even though
individual swaps can move the estimate.

\begin{table}[H]
\centering
\caption{Matched prediction accuracy with different contribution patterns}
\label{tab:scenarios}
\small
\begin{tabular}{lrrrrrr}
\toprule
Error pattern & Gold ATE & Proxy ATE & Gap & RMSE & $R^2$ & Top-two share \\
\midrule
Diffuse & 0.100 & 0.060 & 0.040 & 0.200 & 0.965 & 10\% \\
Concentrated & 0.100 & 0.060 & 0.040 & 0.200 & 0.965 & 100\% \\
Offsetting & 0.100 & 0.100 & 0.000 & 0.200 & 0.965 & 50\% \\
\bottomrule
\end{tabular}


\vspace{0.35em}
\begin{minipage}{0.94\linewidth}
\footnotesize
\textit{Note:} Gold ATE and proxy ATE are differences in means. Gap is gold
minus proxy. Top-two share is the two largest absolute cluster contributions
divided by the sum of absolute contributions. Each row uses the same gold
outcomes. Quantities are deterministic for the generated finite population.
\end{minipage}
\end{table}

The first two patterns share the entire mean swap path, so no average over
validation draws can separate them. Every difference between them is
distributional.

Table~\ref{tab:budgets} prices each validation budget two ways. Tail risk is
the exact share of validation sets whose scaled estimate lands more than
\TailThreshold\ from the gold endpoint. The Normal price is the same
probability computed from a Normal approximation with the correct
finite-population variance of equation~\eqref{eq:htvar}, the summary a
variance-based power analysis would use. The two prices agree only in the
diffuse case, where every draw picks up the same correction and both are zero
at every budget.

Concentrated error breaks the approximation in both directions. At
\BudgetExample\ labeled clusters the Normal price is \ConcentratedNormalTail\
against an exact \ConcentratedTailRisk, because \ConcentratedZeroMove\ of the
draws contain neither high-error cluster and leave the estimate at the
machine-scored value. The exact price is also not monotone in the budget: the
distribution's few atoms cross the threshold at different budgets as the
$G/k$ scaling shrinks, so offsetting risk climbs to \OffsettingTailPeak\ at
\OffsettingTailPeakBudget\ labeled clusters before falling, while the Normal
price declines smoothly and understates the risk where it is largest. A
larger budget can carry more tail risk than a smaller one, and the
variance-based price never shows it. Decision reversals price the same way:
\OffsettingSignFlip\ of single-cluster draws in the offsetting case reverse
the sign of the estimated effect, the analogue of the sign-change criterion in
dropping-data robustness \citep{giordano2026}. Bias is omitted from the table because
it is zero up to numerical precision for every cell. A narrow distribution
does not prove that human labels measure the intended construct. It says only
that the specified validation draw estimates the observed gold-label endpoint
consistently.

\begin{table}[H]
\centering
\caption{Exact design-scaled distributions at selected budgets}
\label{tab:budgets}
\small
\begin{tabular}{lrrrrrr}
\toprule
Error pattern & Gold clusters & Mean & SD & 5th & 95th & Tail risk \\
\midrule
Diffuse & 2 & 0.100 & 0.000 & 0.100 & 0.100 & 0\% \\
 & 5 & 0.100 & 0.000 & 0.100 & 0.100 & 0\% \\
 & 10 & 0.100 & 0.000 & 0.100 & 0.100 & 0\% \\
 & 20 & 0.100 & 0.000 & 0.100 & 0.100 & 0\% \\
\addlinespace
Concentrated & 2 & 0.100 & 0.083 & 0.060 & 0.260 & 19\% \\
 & 5 & 0.100 & 0.048 & 0.060 & 0.220 & 5\% \\
 & 10 & 0.100 & 0.028 & 0.060 & 0.140 & 0\% \\
 & 20 & 0.100 & 0.000 & 0.100 & 0.100 & 0\% \\
\addlinespace
Offsetting & 2 & 0.100 & 0.123 & -0.100 & 0.300 & 35\% \\
 & 5 & 0.100 & 0.071 & 0.020 & 0.180 & 57\% \\
 & 10 & 0.100 & 0.041 & 0.020 & 0.180 & 14\% \\
 & 20 & 0.100 & 0.000 & 0.100 & 0.100 & 0\% \\
\bottomrule
\end{tabular}


\vspace{0.35em}
\begin{minipage}{0.92\linewidth}
\footnotesize
\textit{Note:} Mean, SD, and quantiles describe equation~\eqref{eq:ht} over all
fixed-size cluster samples. Tail risk is
$\Pr(|\widehat\tau_H-\widehat\tau_G|>\TailThreshold)$, computed by exact
enumeration, not Monte Carlo. Normal price is the same probability under a
Normal approximation with the exact finite-population variance of
equation~\eqref{eq:htvar}. Gold ATE equals \GoldATE\ in every scenario.
\end{minipage}
\end{table}

\section{Related subset diagnostics}

A thought experiment locates the swap distribution among its neighbors.
Suppose every observation had a gold label. The measurement layer would
vanish, and the only randomness left would be treatment assignment: over
re-randomization, each cluster's contribution to the estimate enters with a
random sign as the cluster changes arm. The wild cluster bootstrap
approximates exactly this sign-flip process, applying independent mean-zero
multipliers to cluster score contributions \citep{cameron2008}, and the first
break in Section~\ref{sec:breaks} is that layer's variance statement for the
scoring errors. Now remove the labels for all but $k$ clusters. A second
layer of randomness sits on the same cluster contributions: for the linear
contrast,
\[
\widehat\tau_R(S_k)-\E_{S_k}[\widehat\tau_R(S_k)]
=\sum_{g=1}^{G}\left(Z_g-\frac{k}{G}\right)\psi_g,
\]
where $Z_g=\1\{g\in S_k\}$ marks the labeled clusters. The two procedures are
one construction --- recompute the estimate under random reweighting of
cluster contributions --- instantiated with different multiplier laws for
different design layers: signs for assignment, inclusion indicators for
labeling. That is why neither substitutes for the other. Swap quantiles are
not confidence intervals, because valid uncertainty for the experimental
effect must average over the sign layer, and valid uncertainty for the
correction must average over the indicator layer actually used to buy labels.

Dropping-data robustness metrics are the third member of the family: they
apply adversarially chosen zero weights to individual observations and locate
the worst set with a first-order influence approximation, exact enumeration
being infeasible \citep{giordano2026}. The swap construction differs on every
choice --- the perturbed unit is the gold-label set rather than the data, the
weights are drawn by a design rather than chosen by an adversary, and no
approximation is needed because the linear contrast makes every swap effect
exactly additive in the cluster contributions. The adversarial variant is
nevertheless free here: the most adverse $k$-cluster validation draw is the
sum of the $k$ most extreme contributions, read off by sorting the $\psi_g$.

\section{Use and limits}

A validation-swap analysis begins with the same choices required by any
two-phase design. The researcher states the causal estimand, fixes the machine
scoring rule, identifies the unit sampled for human coding, and records every
cluster's inclusion probability. If treatment is assigned by cluster, swapping
individual labels while treating them as independent understates the relevant
variation.

The primary analysis should report a valid correction estimator and interval,
using prediction-powered, design-based supervised learning, or another method
whose assumptions fit the design. The swap distribution is an additional
diagnostic. Before seeing it, the researcher chooses validation budgets and a
substantively meaningful movement threshold. The exhibit then reports
quantiles, tail risk, and any decision reversal probability. Accuracy metrics
remain useful descriptions of the scoring model, but they do not substitute for
this coefficient-specific analysis.

The diagnostic has four limits. First, it cannot reveal the individual errors of
unlabeled observations. Second, a distribution computed inside a convenience
validation sample need not transport outside it. Third, human labels can contain
their own error or define the wrong construct. Fourth, a small number of
validated clusters may leave the distribution too poorly estimated to support
fine subgroup claims. The proper conclusion in that case is that the validation
budget does not resolve the scoring risk.

Machine-scored experiments create two sources of uncertainty: treatment
assignment and gold-label sampling. Randomization-based or cluster-robust
inference addresses the first. Established two-phase estimators address the
second when they correct the machine-scored contrast. The random-swap
distribution makes validation-sample dependence visible by showing how much the
conclusion depends on which clusters were checked.

\bibliographystyle{plainnat}
\bibliography{references}

\end{document}
